\documentclass[10pt,a4paper]{amsart}
\usepackage[margin=19mm]{geometry}
\usepackage{amsmath,amssymb,mathtools}
\usepackage[scr=rsfs,cal=euler]{mathalfa}
\usepackage{xfrac}
\usepackage[unicode,hidelinks,bookmarks=false]{hyperref}
\newcommand{\ie}{i.e.\ }
\newcommand{\cf}{cf.\ }
\newcommand{\resp}{respectively\ }
\newcommand{\Subsection}{\subsection}
\providecommand{\nobreakdash}{\nobreak\hbox{-}}
\setlength{\emergencystretch}{2em}
\multlinegap0pt
\usepackage{bm}
\usepackage{bbm}
\usepackage{dsfont}
\usepackage{xspace}
\usepackage{cancel}
\usepackage{mathtools}
\usepackage{amsthm}
\makeatletter
\@ifundefined{Theorem}{\newtheorem{Theorem}{Theorem}}{}
\makeatother
\DeclareMathOperator{\GL}{GL}
\DeclareMathOperator{\im}{im}
\title[On the Uniqueness of the $G$-Equivariant Spectral Flow]{On the Uniqueness of the $G$-Equivariant Spectral Flow}
\author{Marek Izydorek, Joanna Janczewska, Maciej Starostka, Nils Waterstraat}
\date{Source: arXiv:2508.01061v2 (CC BY 4.0)}
\begin{document}
\maketitle

\noindent\emph{Source and licence.} On the Uniqueness of the $G$-Equivariant Spectral Flow. Version arXiv:2508.01061v2 (CC BY 4.0), distributed under the Creative Commons Attribution 4.0 International licence, as recorded in the arXiv licence metadata for this exact version and shown on the abstract page https://arxiv.org/abs/2508.01061v2. Full rights notice: licenses/arxiv-2508.01061-CC-BY-4.0.txt.\par\smallskip

\noindent\textit{Editorial scope.} Counted unit: the Theorem whose source label is \texttt{thm-cogredpar} in the article (source0.tex lines 258--261, source label \texttt{thm-cogredpar}), delivered together with the article's own complete original proof of it (source0.tex lines 263--365, one \texttt{proof} environment of 103 lines). No mathematics is written, repaired or re-derived: statement and proof are the article's own text line for line. Required context retained uncounted: none. Every definition, display and label of those lines is retained unchanged. Omitted from this package and not redistributed: 1--257, 262--262, 366--851. Nothing inside a retained excerpt is omitted or altered: every retained body is delivered exactly as the article's lines are written, so no omission receipt of any kind accompanies this package. Citations: the retained excerpts carry no citation command at all, so no bibliography entry of the article is transcribed and no reference is redistributed. Reference adaptation: none was needed and none was made; every reference inside the delivered unit resolves inside it, so no unresolved reference or citation is produced. Printed slips kept literal: no printed word, symbol, hypothesis or number of the counted statement or of its proof was altered. Layout and class: the article's own class is \texttt{sigma} with packages including ifpdf, enumitem; the wrapper is the lane's pinned \texttt{amsart} editorial wrapper at 10pt on A4, which restates the theorem-like environments the retained text needs and adds the article's own macro definitions that the retained text needs (\texttt{\textbackslash GL}, \texttt{\textbackslash im}), with extra wrapper packages (bm, bbm, dsfont, xspace, cancel, mathtools). The delivered numbering is the unit's own. Assets: no retained span contains a figure environment, a graphics-inclusion command, a PDF-page inclusion, a table or a TikZ picture; no image file is copied into this package, the package declares no asset dependency and no image file is redistributed with it. Licence: version arXiv:2508.01061v2 of this article is distributed under the Creative Commons Attribution 4.0 International licence, as recorded in the arXiv licence metadata for this exact version and shown on the abstract page https://arxiv.org/abs/2508.01061v2. A full rights notice is delivered at licenses/arxiv-2508.01061-CC-BY-4.0.txt. Characters: every retained excerpt is pure ASCII, so the delivered engine, XeLaTeX with its default Latin Modern text font, reports no missing character.
\begin{Theorem}\label{thm-cogredpar}
Every path $L=\{L_\lambda\}_{\lambda\in I}$ in $\mathcal{FS}^i(H)^G$ has a cogredient parametrix. More precisely, there is a $G$-equivariant symmetry $Q\in\mathcal{FS}^i(H)^G$ as well as paths $M=\{M_\lambda\}_{\lambda\in I}$ in $\GL(H)^G$ and $K=\{K_\lambda\}_{\lambda\in I}$ in $\mathcal{KS}(H)^G$ such that
\[M^\ast_\lambda L_\lambda M_\lambda=Q+K_\lambda,\qquad\lambda\in I.\]
\end{Theorem}

\begin{proof}
We split the proof into three steps.

{\it Step~$1$: The map $\pi_Q$ and local sections.}
We consider for $G$-equivariant symmetries $Q\in\mathcal{FS}^i(H)^G$ the maps
\begin{align}\label{defQ}
\pi_Q \colon \ \GL(H)^G\times\mathcal{KS}(H)^G\rightarrow\mathcal{FS}^i(H)^G,\qquad \pi_Q(M,K)=MQM^\ast+K.
\end{align}
Note that Theorem~\ref{thm-cogredpar} is proved if we can show that there is some $Q$ for which there is a~continuous map $\tilde{L} \colon I\rightarrow \GL(H)^G\times\mathcal{KS}(H)^G$ such that $L_\lambda=\pi_Q\circ\tilde{L}$ for all $\lambda\in I$, i.e., the path $L$ can be lifted to $\GL(H)^G\times\mathcal{KS}(H)^G$ with respect to $\pi_Q$. The right setting to find such a lifting is the theory of fibre bundles and our first aim is to construct local sections for \eqref{defQ}. To be more precise, we are going to show that for any $S\in\mathcal{FS}^i(H)^G$ there is a $G$-equivariant symmetry $Q_S$, an open neighbourhood $\mathcal{U}_S$ of $S$ in $\mathcal{FS}^i(H)^G$ and a map $\sigma_S \colon \mathcal{U}_S\rightarrow\GL(H)^G\times\mathcal{KS}(H)^G$ such that
\[(\pi_{Q_S}\circ\sigma_S)(T)=T\qquad\text{for all} \ T\in\mathcal{U}_S.\]
To construct $\sigma_S$, let $K$ be the orthogonal projection onto the kernel of the selfadjoint Fredholm operator $S$. Then $\ker(S)$ is $G$-invariant and of finite dimension, and thus $K\in\mathcal{KS}(H)^G$ as well~as
\begin{align}\label{V}
V:=S+K\in \GL(H)^G.
\end{align}
To simplify notation, we henceforth let $P_+(V)=\chi_{(0,\infty)}(V)$ and $P_-(V)=\chi_{(-\infty,0)}(V)$ be the orthogonal projections onto the positive and negative spectral subspaces of $V$, which are $G$\nobreakdash-equiv\-ari\-ant. The operator
\begin{align}\label{QS}
Q_S=2P_+(V)-I_H\in\mathcal{FS}^i(H)^G
\end{align}
has a neighbourhood $\tilde{\mathcal{U}}\subset\mathcal{FS}^i(H)^G$ of invertible operators. As before, we denote for $T\in\tilde{\mathcal{U}}$ by $P_+(T)$ and $P_-(T)$ the orthogonal projections onto the positive and negative spectral subspaces. As they continuously depend on $T\in\tilde{\mathcal{U}}$, $Q_S$ has a neighbourhood $\mathcal{U}\subset\tilde{\mathcal{U}}$ such that
\[P_+(T) P_+(Q_S)+P_-(T) P_-(Q_S)\in\GL(H)^G,\qquad T\in \mathcal{U},\]
where it is used that this operator is $I_H\in\GL(H)^G$ for $T=Q_S$. This shows that the restriction of $P_\pm(T)$ to $\im(P_\pm(Q_S))$ is a bijection onto $\im(P_\pm(T))$ for any $T\in\mathcal{U}$. To simplify notation, we~set $H_\pm:=\im(P_\pm(Q_S))$ and consider for $T\in\mathcal{U}$ the map $B(T) \colon H_+\times H_+\rightarrow\mathbb{R}$,
\[B(T)(u,v)=\langle TP_+(T)u,P_+(T)v\rangle,\qquad u,v\in H_+.\]
Then $B(T)$ is positive definite, as well as $G$-invariant as $T$, $P_+(T)$ are $G$-equivariant and $G$ acts orthogonally. Let $T_B$ be the unique selfadjoint operator on $H_+$ such that
\[B(T)(u,v)=\langle T_Bu,v\rangle,\qquad u,v\in H_+,\]
which is usually called the Riesz-representation of $B(T)$. The selfadjoint operator $T_B$, and thus \smash{$S_+(T):=T_B^{\smash{-\frac{1}{2}}}$}, is $G$\nobreakdash-equiv\-ari\-ant as it is unique and $B(T)$ is $G$-invariant. Moreover,
\begin{align*}
\langle u,v\rangle&=\bigl\langle T_B^{-1}T_Bu,v\bigr\rangle=\bigl\langle T_B^{-\frac{1}{2}}T_Bu,T_B^{-\frac{1}{2}}v\bigr\rangle=\bigl\langle T_BT_B^{-\frac{1}{2}}u,T_B^{-\frac{1}{2}}v\bigr\rangle=B(T)\bigl(T_B^{-\frac{1}{2}}u,T_B^{-\frac{1}{2}}v\bigr)\\
&=\bigl\langle TP_+(T)T_B^{-\frac{1}{2}}u,P_+(T)T_B^{-\frac{1}{2}}v\bigr\rangle=\bigl\langle T_B^{-\frac{1}{2}}P_+(T)TP_+(T)T_B^{-\frac{1}{2}}u,v\bigr\rangle\\
&=\langle S_+(T)P_+(T)TP_+(T)S_+(T)u,v\rangle,\qquad u,v\in H_+,
\end{align*}
which implies that
\begin{align}\label{equ-S+}
S_+(T)P_+(T)TP_+(T)S_+(T)=I_{H_+}.
\end{align}
Analogously, there is a map $S_- \colon \mathcal{U}\rightarrow\GL(H_-)^G$ such that for all $T\in\mathcal{U}$
\begin{align}\label{equ-S-}
S_-(T)P_-(T) TP_-(T)S_-(T)=-I_{H_-}.
\end{align}
We now define $S_0 \colon \mathcal{U}\rightarrow\GL(H)^G$ by
\[S_0(T)=P_+(T)S_+(T)P_+(Q_S)-P_-(T)S_-(T)P_-(Q_S),\]
and note that by \eqref{equ-S+} and \eqref{equ-S-}
\begin{align*}
S_0(T)^\ast TS_0(T)=P_+(Q_S)-P_-(Q_S)=Q_S,\qquad T\in\mathcal{U}.
\end{align*}
This shows that
\[\sigma \colon \ \mathcal{U}\rightarrow\GL(H)^G\times\mathcal{KS}(H)^G,\qquad \sigma(T)=((S_0(T)^{-1})^\ast,0),\]
satisfies
\begin{align}\label{equ-picircsigma}
(\pi_{Q_S}\circ\sigma)(T)=\pi_{Q_S}\bigl(\bigl(S_0(T)^{-1}\bigr)^\ast,0\bigr)=\bigl(S_0(T)^{-1}\bigr)^\ast Q_S S_0(T)^{-1}=T,\qquad T\in\mathcal{U}.
\end{align}
Let us recall that $\mathcal{U}$ is a neighbourhood of the symmetry $Q_S$ in \eqref{QS}. We now aim to construct a section $\sigma_V$ in a neighbourhood $\mathcal{U}_V$ of the map $V$ that was defined in \eqref{V}. The product $\GL(H)^G\times\mathcal{KS}(H)^G$ is a closed subgroup of the topological group $\GL(H)\times\mathcal{KS}(H)$, where
\begin{align}\label{groupmult}
(M,K)\cdot \bigl(\tilde{M},\tilde{K}\bigr)=\bigl(M\tilde{M}, K+M\tilde{K}M^\ast\bigr),\qquad (M,K), \bigl(\tilde{M},\tilde{K}\bigr)\in\GL(H)\times\mathcal{KS}(H).
\end{align}
Moreover, there is an action $\tau$ of $\GL(H)^G\times\mathcal{KS}(H)^G$ on $\mathcal{FS}^i(H)^G$ defined by
\[\tau_{h}(L)= MLM^\ast+K,\qquad L\in\mathcal{FS}^i(H),\qquad h=(M,K)\in\GL(H)^G\times\mathcal{KS}(H)^G.\]
For $h:=\bigl(|V|^\frac{1}{2},0\bigr)\in \GL(H)^G\times\mathcal{KS}(H)^G$, we obtain from \eqref{QS}
\[\tau_h(Q_S)=|V|^\frac{1}{2}Q_S|V|^\frac{1}{2}=V\]
and thus see that $\mathcal{U}_V:=\tau_h(\mathcal{U})$ is an open neighbourhood of $V$ in $\mathcal{FS}^i(H)^G$. We set
\[\sigma_V \colon \ \mathcal{U}_V\rightarrow \GL(H)^G\times\mathcal{KS}(H)^G,\qquad \sigma_V(T)=h\cdot \sigma(\tau_{h^{-1}}(T)),\qquad T\in\mathcal{U}_V,\]
where here and henceforth the dot stands for the group multiplication \eqref{groupmult}. To show that $\sigma_V$ is a section of $\pi_{Q_S}$ over $\mathcal{U}_V$, we first note that
\[\pi_{Q_S}(h_1\cdot h_2)=\tau_{h_1}(\pi_{Q_S}(h_2)),\qquad h_1,h_2\in \GL(H)^G\times\mathcal{KS}(H)^G\]
which is a consequence of the definition of $\tau$ and \eqref{groupmult}. Thus, it follows from \eqref{equ-picircsigma} that
\begin{align}
(\pi_{Q_S}\circ\sigma_V)(T)&=\pi_{Q_S}(h\cdot \sigma(\tau_{h^{-1}}(T)))=\tau_h(\pi_{Q_S}(\sigma(\tau_{h^{-1}}(T))))\nonumber\\
&=\tau_h(\tau_{h^{-1}}(T))=T,\qquad T\in\mathcal{U}_V,\label{QSsigmaV}
\end{align}
as claimed. Finally, we set $\tilde{h}=(I_H,-K)\in \GL(H)^G\times\mathcal{KS}(H)^G$, consider the open neighbourhood $\mathcal{U}_S=\tau_{\tilde{h}}(\mathcal{U}_V)$ of $S$ and define
\[\sigma_S\colon \ \mathcal{U}_S\rightarrow \GL(H)^G\times\mathcal{KS}(H)^G,\quad \sigma_S(T)=\tilde{h}\cdot \sigma_V(\tau_{\tilde{h}^{-1}}(T)),\qquad T\in\mathcal{U}_S.\]
The same computation as in \eqref{QSsigmaV} shows that $\pi_{Q_S}\circ\sigma_S=\mathrm{id}|_{\mathcal{U}_S}$, which finishes our first step.

{\it Step $2$: The image of $\pi_Q$.}
The aim of this step of the proof of Theorem~\ref{thm-cogredpar} is to show that for any symmetry $Q$, the image $\pi_Q$ is a union of connected components of $\mathcal{FS}^i(H)^G$. In the argument, we need the auxiliary result that for any symmetries $Q_1$ and $Q_2$
\begin{align}\label{lem33}
\im(\pi_{Q_1})\cap \im(\pi_{Q_2})\neq\varnothing \ \Longrightarrow \ \im(\pi_{Q_1})=\im(\pi_{Q_2})\subset\mathcal{FS}^i(H)^G.
\end{align}
We postpone the proof of \eqref{lem33} and first obtain the actual aim of this step by showing that $\im(\pi_Q)$ is open and closed in $\mathcal{FS}^i(H)^G$. By the first step of the proof, for any $S\in\im(\pi_Q)$ there is a symmetry $Q_S$ and an open neighbourhood $\mathcal{U}_S$ of $S$ in $\mathcal{FS}^i(H)^G$ such that $\mathcal{U}_S\subset\im(\pi_{Q_S})$. Now \eqref{lem33} implies that $\mathcal{U}_S\subset\im(\pi_Q)$ and thus $\im(\pi_Q)$ is open in $\mathcal{FS}^i(H)^G$. Secondly, let $\{S_n\}_{n\in\mathbb{N}}\subset\im(\pi_Q)$ converge to some $S\in\mathcal{FS}^i(H)^G$. Again, by the first step of the proof, there is a symmetry $Q_S$ and a neighbourhood $\mathcal{U}_S$ of $S$ in $\mathcal{FS}^i(H)^G$ such that $\mathcal{U}_S\subset\im(\pi_{Q_S})$. As~$\{S_n\}_{n\in\mathbb{N}}$ converges to $S$, we see that $S_n\in\mathcal{U}_S$ for sufficiently large $n$. By~\eqref{lem33}, this implies that $\mathcal{U}_S\subset\im(\pi_Q)$ and thus $S\in\im(\pi_Q)$, which means that $\im(\pi_Q)$ is closed in $\mathcal{FS}^i(H)^G$.

It remains to show \eqref{lem33}, for which we let $Q_1$, $Q_2$ be two symmetries and assume that $S\in \im(\pi_{Q_1})\cap \im(\pi_{Q_2})$. Then there are $(M_1,K_1), (M_2,K_2)\in\GL(H)^G\times\mathcal{KS}(H)^G$ such that
\[S=\pi_{Q_1}(M_1,K_1)=\pi_{Q_2}(M_2,K_2),\]
and thus
\[S=M_1Q_1M^\ast_1+K_1=M_2Q_2M^\ast_2+K_2,\]
which implies
\[Q_2=M^{-1}_2M_1Q_1M^\ast_1\bigl(M^{-1}_2\bigr)^\ast+M^{-1}_2(K_1-K_2)\bigl(M^{-1}_2\bigr)^\ast.\]
A direct computation shows that
\[\pi_{Q_2}(h)=\pi_{Q_1}\bigl(h\cdot \tilde{h}\bigr),\qquad h\in \GL(H)^G\times\mathcal{KS}(H)^G,\]
where $\tilde{h}=\bigl(M^{-1}_2M_1,M^{-1}_2(K_1-K_2)\bigl(M^{-1}_2\bigr)^\ast\bigr)\in \GL(H)^G\times\mathcal{KS}(H)^G$. This yields $\im(\pi_{Q_2})\subset\im(\pi_{Q_1})$, and proves \eqref{lem33} by swapping $Q_1$ and $Q_2$.

{\it Step $3$: Fibre bundles and end of the proof.}
Henceforth, we set $B_Q=\im(\pi_Q)$ for any given symmetry $Q\in\mathcal{FS}^i(H)^G$. Moreover, let us note that if $S=Q$ in Step 1, then it follows from \eqref{V} and \eqref{QS} that $Q_Q=Q$, and thus for every symmetry $Q\in\mathcal{FS}^i(H)^G$ there is some $S\in\mathcal{FS}^i(H)^G$ such that $Q_S=Q$. Therefore, it is not necessary to specify $S$ in our notation any longer.

We now first note that the map $\pi_Q \colon \GL(H)^G\times\mathcal{KS}(H)^G\rightarrow B_Q$ is the projection of a locally trivial fibre-bundle with fibre given by the isotropy group of $Q\in\mathcal{FS}^i(H)^G$. Indeed, by the first step of the proof, there is a local section $\sigma$ of $\pi_Q$ on an open subset $\mathcal{U}_Q\subset B_Q$. Now a local trivialisation over $\mathcal{U}_Q$ is given by
\[\eta \colon \ \mathcal{U}_Q\times\pi^{-1}_Q(Q)\rightarrow\pi^{-1}_Q(\mathcal{U}_Q),\qquad \eta(S,h)=\sigma(S)\cdot h,\]
and it can be shifted to any point $T\in B_Q$ in the following usual way. As $\im(\pi_Q)=B_Q$ by definition, there is some element $\tilde{h}\in \GL(H)^G\times\mathcal{KS}(H)^G$ such that $\pi_Q(\tilde{h})=\tau_{\tilde{h}}(Q)=T$. Now $\mathcal{U}:=\tau_{\tilde{h}}(\mathcal{U}_Q)$ is a neighbourhood of $T$ and $\tau_{\tilde{h}} \colon \mathcal{U}_Q\rightarrow\mathcal{U}$ is a homeomorphism, which yields a local trivialisation over $\mathcal{U}$ by
\[\eta' \colon \ \mathcal{U}\times\pi^{-1}_Q(Q)\rightarrow\pi^{-1}_Q(\mathcal{U}),\qquad \eta'(S,h)=\tilde{h}\cdot\sigma(\tau_{\tilde{h}^{-1}}(S))\cdot h.\]
Now we finally can prove Theorem~\ref{thm-cogredpar} by standard fibre bundle theory. If $L=\{L_\lambda\}_{\lambda\in I}$ is a~path in $\mathcal{FS}^i(H)^G$, then $L(I)$ is contained in a path component of $\mathcal{FS}^i(H)^G$, which we henceforth call~$C$. Let $S\in C$ be arbitrary, $Q$ the associated symmetry as in the first step and note that $C\subset B_{Q}$ by the second step. Now we consider the pullback $(E,I,\pi)$ of the fibre bundle
\begin{align}\label{piQ_S}
\pi_{Q} \colon \ \GL(H)^G\times\mathcal{KS}(H)^G\rightarrow B_{Q}
\end{align}
by $L$. This bundle has the set
\[E=\bigl\{(\lambda,h)\in I\times \bigl(\GL(H)^G\times\mathcal{KS}(H)^G\bigr) \mid L_\lambda=\pi_Q(h)\bigr\}\]
as total space and the bundle projection $\pi$ is the restriction of the projection onto the first component. As $I$ is contractible, $(E,I,\pi)$ is trivial and consequently has a globally defined section. Now the projection onto the second component $I\times\bigl(\GL(H)^G\times\mathcal{KS}(H)^G\bigr)\rightarrow\GL(H)^G\times\mathcal{KS}(H)^G$ is a bundle map from $E$ to the total space of \eqref{piQ_S}. By composing with this map, sections of $(E,I,\pi)$ yield liftings of $L$. Thus, we~have found the desired map $\tilde{L} \colon I\rightarrow\GL(H)^G\times\mathcal{KS}(H)^G$ such that $L_\lambda=\pi_Q\circ\tilde{L}_\lambda$ for all $\lambda\in I$ and so Theorem~\ref{thm-cogredpar} is proved.
\end{proof}

\end{document}
